% SOURCE_STATEMENT: "What are the essential ingredients of a non-formal deformation quantization, including the associated Hilbert-algebra or von Neumann-algebra structure?"
% Preparation and final self-audit: CONTRIBUTING.md prompts C and D; source/audit artifacts retained in the local physics_statements_2026_10_11 batch.
% !TeX program = lualatex
\documentclass[11pt]{article}
\usepackage[margin=25mm]{geometry}
\usepackage{fontspec}
\setmainfont{Latin Modern Roman}
\usepackage{amsmath,amssymb,mathtools}
\usepackage{unicode-math}
\setmathfont{Latin Modern Math}
\usepackage{xurl,hyperref}
\hypersetup{colorlinks=true,urlcolor=blue}
\setcounter{secnumdepth}{0}
\setlength{\parindent}{0pt}
\setlength{\parskip}{5pt}
\setlength{\emergencystretch}{3em}
\newcommand{\sref}[2]{\hyperlink{source:#1:#2}{[S#2]}}
\newcommand{\cataloguescope}[1]{\paragraph{Recorded target.}#1}
\newcommand{\R}{\mathbb R}
\newcommand{\C}{\mathbb C}
\newcommand{\N}{\mathbb N}
\newcommand{\Z}{\mathbb Z}
\newcommand{\Q}{\mathbb Q}
\newcommand{\F}{\mathbb F}
\newcommand{\eps}{\varepsilon}
\begin{document}
\section{Analytic Structures for Nonformal Deformation Quantization}\label{statement}

\subsection{Definitions and mathematical statement}

\paragraph{Shared notation.}$\mathbb N=\{1,2,\ldots\}$; the standard number fields have their usual meanings. All additional conventions are stated locally.


\paragraph{Definitions and conventions.}
A formal deformation quantization replaces multiplication by $f\star g=fg+\sum_{j\geq1}\hbar^j C_j(f,g)$ in formal power series, with the first-order commutator encoding the Poisson bracket. Non-formal quantization instead requires actual parameter values and analytic control of products or representations. A Hilbert algebra combines an involutive algebra with a compatible inner product; a von Neumann algebra is a weakly closed operator algebra. These analytic structures require domains, positivity and completion conventions.

\paragraph{Target retained from the canonical source catalogue.}
\paragraph{Catalogue wording retained for traceability.}
What are the essential ingredients of a non-formal deformation quantization, including the associated Hilbert-algebra or von Neumann-algebra structure?

\paragraph{Source scope and qualification.}
This is a foundations programme, not a claim that every Poisson manifold admits a specified strict quantization. The desired notion of non-formal quantization must be made explicit in any proposed answer.
\subsection{Short English statement}
Identify the analytic ingredients that turn formal deformation quantization into an actual Hilbert- or operator-algebra theory.
\subsection{Sources}
\begin{itemize}
\item[\hypertarget{source:statement:1}{S1}] ULAM.AI and the UnsolvedMath Contributors, UnsolvedMath, catalogue entry. Catalogue curation under CC BY 4.0; definitions and English exposition added locally.
\url{https://huggingface.co/datasets/ulamai/UnsolvedMath}.
\item[\hypertarget{source:statement:2}{S2}] Mini-Workshop: Deformation quantization: between formal to strict (2015), Oberwolfach Reports, DOI 10.4171/owr/2015/11, pp. 571--595. Primary report accessed 2026-10-11; relevant PDF pages 10.
\url{https://doi.org/10.4171/owr/2015/11}.
\item[\hypertarget{source:statement:3}{S3}] Rieffel, Quantization and C*-algebras (catalogue-listed context; not independently consulted in this preparation).
\url{https://math.berkeley.edu/~rieffel/papers/quantization.pdf}.
\end{itemize}
\label{end:statement}
\end{document}
